\section{集合的基本概念及其运算 III}

\subsection{有序偶与笛卡儿乘积}

\begin{definition}[有序偶]
	将对象 $x, y$ 按指定顺序组成的序列称为\textbf{有序偶}，记作 $\langle x, y \rangle$。可以用集合将其定义为
	$$
	\langle x, y \rangle = \big\{\{x\}, \{x, y\}\big\}。
	$$
\end{definition}

有序偶满足唯一性：
$$
\langle u, v \rangle = \langle x, y \rangle
\iff u = x \land v = y。
$$
因此有序偶不同于二元集；一般而言，$\langle x, y \rangle \neq \langle y, x \rangle$，而 $\{x, y\} = \{y, x\}$。

\begin{definition}[笛卡儿乘积]
	集合 $A$ 与 $B$ 的\textbf{笛卡儿乘积}定义为
	$$
	A \times B = \{\langle x, y \rangle \mid x \in A \land y \in B\}。
	$$
\end{definition}

笛卡儿乘积一般不满足交换律。它具有以下基本性质：
$$
\begin{aligned}
	A \times B = \varnothing
	& \iff A = \varnothing \lor B = \varnothing, \\
	\abs{A \times B}
	& = \abs{A}\abs{B} && \text{$A, B$ 为有穷集}, \\
	A \times B \subseteq C \times D
	& \iff A \subseteq C \land B \subseteq D
	&& \text{$A, B, C, D$ 非空}。
\end{aligned}
$$

笛卡儿乘积对交、并、差满足分配律，例如
$$
\begin{aligned}
	A \times (B \cup C) & = (A \times B) \cup (A \times C), \\
	A \times (B \cap C) & = (A \times B) \cap (A \times C), \\
	A \times (B \setminus C) & = (A \times B) \setminus (A \times C)。
\end{aligned}
$$
对第一分量也有相应的等式。

更一般地，定义
$$
A_1 \times A_2 \times \dots \times A_n
= \{\langle x_1, x_2, \dots, x_n \rangle \mid x_i \in A_i,\ 1 \le i \le n\}，
$$
并将 $A \times A \times \dots \times A$ 简记为 $A^n$。

可以使用归纳定义的方法定义多元组：
$$
\langle x_1, x_2, \dots, x_n \rangle := \langle \langle x_1, x_2, \dots, x_{n - 1} \rangle, x_n \rangle
$$
